% English slide deck: beamer, 16:9, XeLaTeX. No bibliography — cite sources
% in a short "Source:" line or footnote on the frame that uses them.
% Structure: title -> outline -> motivation -> method (equation) -> results
% (booktabs table) -> take-home messages -> closing frame.
% Rule of thumb: one message per frame; frame count is not talk length.
\documentclass[aspectratio=169]{beamer}

\usetheme{Boadilla}
\setbeamertemplate{navigation symbols}{}
\usepackage{amsmath,amssymb}
\usepackage{booktabs}
\usepackage{graphicx}

\title{Talk Title}
\subtitle{Optional Subtitle}
\author{Speaker Name}
\institute{Affiliation}
% Name the event and date; \today would print the date this version was
% first built (the build clock is pinned per snapshot for reproducibility).
\date{Event Name, Month Year}

\begin{document}

\begin{frame}
  \titlepage
\end{frame}

\begin{frame}{Outline}
  \tableofcontents
\end{frame}

\section{Motivation}
\begin{frame}{Motivation}
  % The problem, why this audience should care, and the one idea to remember.
  \begin{itemize}
    \item The problem this talk addresses
    \item Why it matters to this audience
    \item The one idea to remember
  \end{itemize}
\end{frame}

\section{Method}
\begin{frame}{Key Idea}
  % Replace with the central equation; define every symbol on the slide.
  \begin{equation}
    \hat{\theta} = \operatorname*{arg\,min}_{\theta}
      \sum_{i=1}^{n} \ell\bigl(f_{\theta}(x_i),\, y_i\bigr)
  \end{equation}
  \begin{itemize}
    \item $f_{\theta}$: model with parameters $\theta$
    \item $\ell$: loss comparing prediction and target
  \end{itemize}
\end{frame}

\section{Results}
\begin{frame}{Results}
  % Report measured numbers only; say where they come from.
  \begin{table}
    \centering
    \begin{tabular}{lcc}
      \toprule
      Method   & Metric A & Metric B \\
      \midrule
      Baseline & --       & --       \\
      Proposed & --       & --       \\
      \bottomrule
    \end{tabular}
    \caption{Placeholder header; fill in measured results.}
  \end{table}
\end{frame}

\section{Conclusion}
\begin{frame}{Take-home Messages}
  \begin{enumerate}
    \item First message
    \item Second message
    \item Limitations and next steps
  \end{enumerate}
\end{frame}

\begin{frame}
  \centering
  \Large Thank you --- questions?
\end{frame}

\end{document}
